\documentclass{article}
\usepackage{iftex}
\ifXeTeX
  \usepackage{fontspec}
  \usepackage{xeCJK}
  \setmonofont[Path=/Users/yupyom/.src2tex/fonts/hackgen/, BoldFont=HackGen-Bold.ttf]{HackGen-Regular.ttf}
  \setCJKmonofont[Path=/Users/yupyom/.src2tex/fonts/hackgen/, BoldFont=HackGen-Bold.ttf]{HackGen-Regular.ttf}
  \setCJKmainfont[Path=/usr/local/texlive/2026/texmf-dist/fonts/opentype/public/haranoaji/, Extension=.otf, BoldFont=HaranoAjiMincho-Bold.otf]{HaranoAjiMincho-Regular}
  \XeTeXlinebreaklocale ""
  \xeCJKsetup{CJKglue={},CJKecglue={}}
\fi
\ifLuaTeX
  \usepackage{luatexja-fontspec}
  \setmonofont[Path=/Users/yupyom/.src2tex/fonts/hackgen/, BoldFont=HackGen-Bold.ttf]{HackGen-Regular.ttf}
  \setmainjfont{HaranoAjiMincho-Regular.otf}
  \ltjsetparameter{kanjiskip=0pt, xkanjiskip=0pt}
\fi
\usepackage[a4paper, margin=2cm]{geometry}
\usepackage{graphicx}
\newdimen\charwd
{\tt\global\setbox0=\hbox{x}\global\charwd=\wd0}
\frenchspacing
\usepackage{fancyhdr}
\pagestyle{fancy}
\renewcommand{\headrulewidth}{0pt}
\fancyhf{}
\fancyhead[R]{\rm src2tex-go version 2.23}
% plain TeX compatibility
\makeatletter
\providecommand{\eqalign}[1]{\vcenter{\openup1\jot\m@th
  \ialign{\strut\hfil$\displaystyle{##}$&$\displaystyle{{}##}$\hfil\crcr#1\crcr}}}
\makeatother
\fancyfoot[R]{\rm\hfill samples{\tt /}hanoi.py\qquad page \thepage}
\begin{document}

\ifx\sevenrm\undefined
  \font\sevenrm=cmr7 scaled \magstep0
\fi

\newread\MyStyle
\openin\MyStyle=src2latex.s2t
\ifeof\MyStyle
  \closein\MyStyle
\else
  \input src2latex.s2t
  \closein\MyStyle
\fi

\def\mc{\relax}
\def\gt{\relax}
\def\sc{\scshape}

\tt\mc 

\noindent
\rm\mc {\tt\#}\  \hrulefill  \rm\mc 

\noindent
\mbox{}{\tt\#}
\tt\mc 

\noindent
\mbox{}\rm\mc {\tt\#}\  {}% beginning of TeX mode  \rm\mc 

\noindent
\mbox{}{\tt\#}
\tt\mc 

\noindent
\mbox{}\rm\mc {\tt\#}\  {}\centerline{\bf Towers of Hanoi (Python)}  \rm\mc 

\noindent
\mbox{}{\tt\#}
\tt\mc 

\noindent
\mbox{}\rm\mc {\tt\#}\  {}  \rm\mc 

\noindent
\mbox{}{\tt\#}\  {}\hspace{\leftmargini}\noindent  \rm\mc 

\noindent
\mbox{}{\tt\#}\  {}\hspace{\leftmargini}This program gives an answer to the following famous problem (towers of  \rm\mc 

\noindent
\mbox{}{\tt\#}\  {}\hspace{\leftmargini}Hanoi).  \rm\mc 

\noindent
\mbox{}{\tt\#}\  {}\hspace{\leftmargini}There is a legend that when one of the temples in Hanoi was constructed,  \rm\mc 

\noindent
\mbox{}{\tt\#}\  {}\hspace{\leftmargini}three poles were erected and a tower consisting of 64 golden discs was  \rm\mc 

\noindent
\mbox{}{\tt\#}\  {}\hspace{\leftmargini}arranged on one pole, their sizes decreasing regularly from bottom to top.  \rm\mc 

\noindent
\mbox{}{\tt\#}\  {}\hspace{\leftmargini}The monks were to move the tower of discs to the opposite pole, moving  \rm\mc 

\noindent
\mbox{}{\tt\#}\  {}\hspace{\leftmargini}only one at a time, and never putting any size disc above a smaller one.  \rm\mc 

\noindent
\mbox{}{\tt\#}\  {}\hspace{\leftmargini}The job was to be done in the minimum numbers of moves. What strategy for  \rm\mc 

\noindent
\mbox{}{\tt\#}\  {}\hspace{\leftmargini}moving discs will accomplish this optimum transfer?  \rm\mc 

\noindent
\mbox{}{\tt\#}\  {}  \rm\mc 

\noindent
\mbox{}{\tt\#}
\tt\mc 

\noindent
\mbox{}\rm\mc {\tt\#}\  {}% end of TeX mode  \rm\mc 

\noindent
\mbox{}{\tt\#}
\tt\mc 

\noindent
\mbox{}\rm\mc {\tt\#}\  \hrulefill  \rm\mc 

\noindent
\mbox{}\tt\mc \hfill

\noindent
\mbox{}\hfill

\noindent
\mbox{}\rm\mc {\tt\#}\  \hrulefill\ hanoi.py \ \hrulefill \rm\mc 

\noindent
\mbox{}\tt\mc \hfill

\noindent
\mbox{}\hfill

\noindent
\mbox{}ARRAY{\tt\mc \ }={\tt\mc \ }8{\tt\mc \ \ \ \ \ \ \ \ \ \ \ \ \ \ \ \ \ \ \ \ \ \ \ \ \ \ \ \ \ \ \ }\rm\mc {\tt\#}\  {}disc の数 \hfill  \rm\mc 

\noindent
\mbox{}\tt\mc \hfill

\noindent
\mbox{}disc{\tt\mc \ }={\tt\mc \ }[[0]{\tt *}ARRAY{\tt\mc \ }for{\tt\mc \ }{\tt\_}{\tt\mc \ }in{\tt\mc \ }range(3)]

\noindent
\mbox{}{\tt\mc \ \ \ \ \ \ \ \ \ \ \ \ \ \ \ \ \ \ \ \ \ \ \ \ \ \ \ \ \ \ \ \ \ \ \ \ \ \ \ \ }\rm\mc {\tt\#}\  {}disc に関するデータの置き場所 \hfill  \rm\mc 

\noindent
\mbox{}\tt\mc \hfill

\noindent
\mbox{}def{\tt\mc \ }init{\tt\_}array():{\tt\mc \ \ \ \ \ \ \ \ \ \ \ \ \ \ \ \ \ \ \ \ \ \ \ }\rm\mc {\tt\#}\  {}disc に関するデータの初期化 \hfill  \rm\mc 

\noindent
\mbox{}\tt\mc {\tt\mc \ \ \ \ }for{\tt\mc \ }j{\tt\mc \ }in{\tt\mc \ }range(ARRAY):

\noindent
\mbox{}{\tt\mc \ \ \ \ \ \ \ \ }disc[0][j]{\tt\mc \ }={\tt\mc \ }ARRAY{\tt\mc \ }{\tt -}{\tt\mc \ }j

\noindent
\mbox{}{\tt\mc \ \ \ \ \ \ \ \ }disc[1][j]{\tt\mc \ }={\tt\mc \ }0

\noindent
\mbox{}{\tt\mc \ \ \ \ \ \ \ \ }disc[2][j]{\tt\mc \ }={\tt\mc \ }0

\noindent
\mbox{}\hfill

\noindent
\mbox{}counter{\tt\mc \ }={\tt\mc \ }0{\tt\mc \ \ \ \ \ \ \ \ \ \ \ \ \ \ \ \ \ \ \ \ \ \ \ \ \ \ \ \ \ }\rm\mc {\tt\#}\  {}移動回数カウンタ \hfill  \rm\mc 

\noindent
\mbox{}\tt\mc \hfill

\noindent
\mbox{}def{\tt\mc \ }print{\tt\_}result():{\tt\mc \ \ \ \ \ \ \ \ \ \ \ \ \ \ \ \ \ \ \ \ \ }\rm\mc {\tt\#}\  {}結果の表示 \hfill  \rm\mc 

\noindent
\mbox{}\tt\mc {\tt\mc \ \ \ \ }global{\tt\mc \ }counter

\noindent
\mbox{}{\tt\mc \ \ \ \ }counter{\tt\mc \ }+={\tt\mc \ }1

\noindent
\mbox{}{\tt\mc \ \ \ \ }print(f{\tt "}{\tt -}{\tt -}{\tt -}{\tt\#}{\tt\char'173}counter{\tt\char'175}{\tt -}{\tt -}{\tt -}{\tt "})

\noindent
\mbox{}{\tt\mc \ \ \ \ }for{\tt\mc \ }i{\tt\mc \ }in{\tt\mc \ }range(3):

\noindent
\mbox{}{\tt\mc \ \ \ \ \ \ \ \ }print(f{\tt "}[{\tt\char'173}i{\tt\char'175}]{\tt "},{\tt\mc \ }end={\tt "}{\tt\mc \ }{\tt "})

\noindent
\mbox{}{\tt\mc \ \ \ \ \ \ \ \ }for{\tt\mc \ }j{\tt\mc \ }in{\tt\mc \ }range(ARRAY):

\noindent
\mbox{}{\tt\mc \ \ \ \ \ \ \ \ \ \ \ \ }if{\tt\mc \ }disc[i][j]{\tt\mc \ }!={\tt\mc \ }0:

\noindent
\mbox{}{\tt\mc \ \ \ \ \ \ \ \ \ \ \ \ \ \ \ \ }print(disc[i][j],{\tt\mc \ }end={\tt "}{\tt\mc \ }{\tt "})

\noindent
\mbox{}{\tt\mc \ \ \ \ \ \ \ \ \ \ \ \ }else:

\noindent
\mbox{}{\tt\mc \ \ \ \ \ \ \ \ \ \ \ \ \ \ \ \ }break

\noindent
\mbox{}{\tt\mc \ \ \ \ \ \ \ \ }print()

\noindent
\mbox{}\hfill

\noindent
\mbox{}ptr{\tt\mc \ }={\tt\mc \ }[0,{\tt\mc \ }0,{\tt\mc \ }0]{\tt\mc \ \ \ \ \ \ \ \ \ \ \ \ \ \ \ \ \ \ \ \ \ \ \ \ \ }\rm\mc {\tt\#}\  {}disc 移動用ポインタ \hfill  \rm\mc 

\noindent
\mbox{}\tt\mc \hfill

\noindent
\mbox{}def{\tt\mc \ }move{\tt\_}one{\tt\_}disc(i,{\tt\mc \ }j):{\tt\mc \ \ \ \ \ \ \ \ \ \ \ \ \ \ \ \ }\rm\mc {\tt\#}\  {}1枚の disc を pole $i$ から \hfill  \rm\mc 

\noindent
\mbox{}\tt\mc {\tt\mc \ \ \ \ \ \ \ \ \ \ \ \ \ \ \ \ \ \ \ \ \ \ \ \ \ \ \ \ \ \ \ \ \ \ \ \ \ \ \ \ }\rm\mc {\tt\#}\  {}pole $j$ に移動する \hfill  \rm\mc 

\noindent
\mbox{}\tt\mc {\tt\mc \ \ \ \ }ptr[i]{\tt\mc \ }{\tt -}={\tt\mc \ }1

\noindent
\mbox{}{\tt\mc \ \ \ \ }disc[j][ptr[j]]{\tt\mc \ }={\tt\mc \ }disc[i][ptr[i]]

\noindent
\mbox{}{\tt\mc \ \ \ \ }ptr[j]{\tt\mc \ }+={\tt\mc \ }1

\noindent
\mbox{}{\tt\mc \ \ \ \ }disc[i][ptr[i]]{\tt\mc \ }={\tt\mc \ }0

\noindent
\mbox{}\hfill

\noindent
\mbox{}def{\tt\mc \ }move{\tt\_}discs(n,{\tt\mc \ }i,{\tt\mc \ }j,{\tt\mc \ }k):{\tt\mc \ \ \ \ \ \ \ \ \ \ \ \ \ }\rm\mc {\tt\#}\  {}上から $n$ 枚目までの disc を \hfill  \rm\mc 

\noindent
\mbox{}\tt\mc {\tt\mc \ \ \ \ \ \ \ \ \ \ \ \ \ \ \ \ \ \ \ \ \ \ \ \ \ \ \ \ \ \ \ \ \ \ \ \ \ \ \ \ }\rm\mc {\tt\#}\  {}pole $i$ から pole $j$ に \hfill  \rm\mc 

\noindent
\mbox{}\tt\mc {\tt\mc \ \ \ \ \ \ \ \ \ \ \ \ \ \ \ \ \ \ \ \ \ \ \ \ \ \ \ \ \ \ \ \ \ \ \ \ \ \ \ \ }\rm\mc {\tt\#}\  {}pole $k$ を経由して移動する \hfill  \rm\mc 

\noindent
\mbox{}\tt\mc {\tt\mc \ \ \ \ }if{\tt\mc \ }n{\tt\mc \ }{\tt >}={\tt\mc \ }1:

\noindent
\mbox{}{\tt\mc \ \ \ \ \ \ \ \ }move{\tt\_}discs(n{\tt -}1,{\tt\mc \ }i,{\tt\mc \ }k,{\tt\mc \ }j)

\noindent
\mbox{}{\tt\mc \ \ \ \ \ \ \ \ \ \ \ \ \ \ \ \ \ \ \ \ \ \ \ \ \ \ \ \ \ \ \ \ \ \ \ \ \ \ \ \ }\rm\mc {\tt\#}\  {}関数 {\tt move\_discs()} の中で \hfill  \rm\mc 

\noindent
\mbox{}\tt\mc {\tt\mc \ \ \ \ \ \ \ \ \ \ \ \ \ \ \ \ \ \ \ \ \ \ \ \ \ \ \ \ \ \ \ \ \ \ \ \ \ \ \ \ }\rm\mc {\tt\#}\  {}さらに自分自身が使われている。 \hfill  \rm\mc 

\noindent
\mbox{}\tt\mc {\tt\mc \ \ \ \ \ \ \ \ }move{\tt\_}one{\tt\_}disc(i,{\tt\mc \ }j){\tt\mc \ \ \ \ \ \ \ \ \ \ \ \ \ }\rm\mc {\tt\#}\  {}このような手法は「再帰的呼び \hfill  \rm\mc 

\noindent
\mbox{}\tt\mc {\tt\mc \ \ \ \ \ \ \ \ }print{\tt\_}result(){\tt\mc \ \ \ \ \ \ \ \ \ \ \ \ \ \ \ \ \ \ }\rm\mc {\tt\#}\  {}だし」といわれる。 \hfill  \rm\mc 

\noindent
\mbox{}\tt\mc {\tt\mc \ \ \ \ \ \ \ \ }move{\tt\_}discs(n{\tt -}1,{\tt\mc \ }k,{\tt\mc \ }j,{\tt\mc \ }i)

\noindent
\mbox{}\hfill

\noindent
\mbox{}\rm\mc {\tt\#}\  {}\par\begin{center}\includegraphics[scale=0.3]{hanoi1}\quad\includegraphics[scale=0.3]{hanoi2}\end{center}  \rm\mc 

\noindent
\mbox{}{\tt\#}
\tt\mc 

\noindent
\mbox{}\rm\mc {\tt\#}\  {}\noindent  \rm\mc 

\noindent
\mbox{}{\tt\#}\  {}たとえば、関数 {\tt move\_discs(4, 0, 1, 2)} を呼び出すことは、  \rm\mc 

\noindent
\mbox{}{\tt\#}\  {}上図のような操作をすることに対応する。\hfill  \rm\mc 

\noindent
\mbox{}\tt\mc \hfill

\noindent
\mbox{}\hfill

\noindent
\mbox{}if{\tt\mc \ }{\tt\_}{\tt\_}name{\tt\_}{\tt\_}{\tt\mc \ }=={\tt\mc \ }{\tt "}{\tt\_}{\tt\_}main{\tt\_}{\tt\_}{\tt "}:

\noindent
\mbox{}{\tt\mc \ \ \ \ }ptr[0]{\tt\mc \ }={\tt\mc \ }ARRAY

\noindent
\mbox{}{\tt\mc \ \ \ \ }ptr[1]{\tt\mc \ }={\tt\mc \ }0

\noindent
\mbox{}{\tt\mc \ \ \ \ }ptr[2]{\tt\mc \ }={\tt\mc \ }0

\noindent
\mbox{}\hfill

\noindent
\mbox{}{\tt\mc \ \ \ \ }init{\tt\_}array()

\noindent
\mbox{}{\tt\mc \ \ \ \ }move{\tt\_}discs(ARRAY,{\tt\mc \ }0,{\tt\mc \ }1,{\tt\mc \ }2){\tt\mc \ \ \ \ \ \ \ \ \ \ }\rm\mc {\tt\#}\  {}{\tt ARRAY} 枚の disc を \hfill  \rm\mc 

\noindent
\mbox{}\tt\mc {\tt\mc \ \ \ \ \ \ \ \ \ \ \ \ \ \ \ \ \ \ \ \ \ \ \ \ \ \ \ \ \ \ \ \ \ \ \ \ \ \ \ \ }\rm\mc {\tt\#}\  {}pole 0 から pole 1 に \hfill  \rm\mc 

\noindent
\mbox{}\tt\mc {\tt\mc \ \ \ \ \ \ \ \ \ \ \ \ \ \ \ \ \ \ \ \ \ \ \ \ \ \ \ \ \ \ \ \ \ \ \ \ \ \ \ \ }\rm\mc {\tt\#}\  {}pole 2 を経由して移動する \hfill  \rm\mc 

\noindent
\mbox{}

\rm\mc

\end{document}
